%%%%%%%%%%  scan idx 33  |  folio 20  %%%%%%%%%%
\begin{center}
{\Large \underline{\underline{CLASSES DE FAISCEAUX ET THÉORÈME DE RIEMANN--ROCH}}}\par
\vspace{0.5em}
\underline{par A. GROTHENDIECK}\footnote{Ceci est la reproduction textuelle du rapport cité dans l'Introduction du présent Séminaire. Les notes de bas de page indiquées par des signes (*), (**) ont été rajoutées en Octobre 1967.}
\end{center}

\subsection*{CHAP. I -- $\lambda$-\uline{Anneaux} (\uline{préliminaires formels})}
\begin{longtable}{@{}p{0.74\linewidth}r@{}}
1. Définitions & 1\\
2. Exemples & 4\\
3. Le $\lambda$-anneau défini par un anneau gradué & 8\\
4. Les opérations $\lambda^p(N,x)$ & 13\\
\end{longtable}

\subsection*{CHAP. II -- \uline{Classes de faisceaux algébriques cohérents et classes de Chern}\hfill 18}
\begin{longtable}{@{}p{0.74\linewidth}r@{}}
1. La théorie de Chow & 18\\
2. Définition des classes de Chern des faisceaux algébriques cohérents & 22\\
3. Généralités fonctorielles sur $K(X)$ & 27\\
4. Quelques résultats techniques & 34\\
5. Définition faisceautique des classes de Chern. Application à l'étude des morphismes d'injection & 43\\
6. Le théorème de Riemann--Roch & 49
\end{longtable}

\begin{center}
Démonstration du théorème de Riemann--Roch (1er Novembre 1957).
\end{center}

\clearpage

